\section{Review of Existing Knowledge}\label{sec:chapter-2}

\subsection{Content-based and collaborative recommendation}

Content-based recommendation matches descriptions of items to a representation of a user's interests or current context. For Holistic Mind, exercise descriptions, categories, intended states, and support goals provide this content. The approach can produce suggestions before an account has accumulated interaction history. Its weakness is dependence on the quality of metadata and the chosen similarity representation. Missing restrictions or ambiguous wording remain problematic even when a numerical similarity score is high.

Collaborative recommendation uses patterns in interaction data to estimate which items may suit a user. Holistic Mind includes a neighbour-based component, but enables it only when sufficient overlapping interactions and neighbours exist. This is appropriate because sparse histories cannot support a dependable comparison. Completion, repetition, saving, and explicit helpfulness feedback also represent different behaviours. Treating every event as equally positive would hide those distinctions. The implementation consequently assigns different values to events and prioritises uncomfortable feedback as a negative signal.

\subsection{Hybrid approaches}

\citet{burke2002} surveys ways to combine recommendation techniques and explains why hybrid systems can address weaknesses of individual methods. Holistic Mind applies that principle through eligibility filtering followed by weighted rules, semantic relevance, and conditional collaborative scoring. The rationale is complementary evidence: structured rules support predictable alignment, embeddings capture descriptive similarity, and interactions may support later adaptation. This rationale does not prove that any particular weighting is optimal; weight selection remains a separate empirical problem.

A key distinction is between eligibility and ranking. Eligibility decides whether an exercise may appear at all under the application's known restrictions. Ranking orders the remaining candidates. If restrictions were only soft penalties, sufficiently strong similarity could override them. The implemented hard exclusions avoid that failure for declared concerns. Nevertheless, they cannot address a restriction that was never collected or encoded in metadata. The system therefore needs both transparent data collection and conservative interpretation of its output.

\subsection{Semantic representations}

\citet{reimers2019} describe Sentence-BERT representations that support efficient comparison of sentences using cosine similarity. Holistic Mind uses the pretrained all-MiniLM-L6-v2 model through ONNX inference. Its model card describes a 384-dimensional representation suited to sentence and short-paragraph tasks \citep{minilm}. This supports a practical comparison between check-in descriptions and exercise documents without requiring a generative model to compose advice. The project does not fine-tune MiniLM or claim a newly trained language model.

Semantic similarity remains narrower than understanding a person's needs. A vector representation may associate related words without distinguishing whether an activity is welcome, uncomfortable, or medically inappropriate. The current engine consequently limits semantic scoring to eligible candidates and combines it with rules. Model failure is also observable: the service can identify a lexical fallback backend, while strict ONNX evaluation aborts rather than silently producing results labelled as semantic inference. This distinction makes benchmark interpretation more defensible.

\subsection{Privacy and evaluation gaps}

\citet{morris2023} demonstrate that text can be recovered from embeddings under studied attack conditions. Their work challenges the assumption that converting a journal to a vector makes its meaning private. Holistic Mind therefore uploads neither journal text nor device-generated journal embeddings for personalisation. This is a specific application design decision, not a claim that every embedding system has identical leakage. It also creates an explicit trade-off: the recommender has less reflective context available.

The reference dissertations provide useful examples of connecting architecture, implementation evidence, metrics, and critical evaluation. Rai's travel recommender and Fullel's Nepali voice assistant address different domains, datasets, and users. Their performance figures cannot validate Holistic Mind. The gap addressed here is the combination of a daily wellness workflow, declared comfort controls, and a journal privacy boundary within a small recommendation application. Clinical effectiveness, longitudinal engagement, and representative usability remain future research questions rather than established contributions.

\subsection{Implications for this project}

The reviewed ideas suggest three practical design principles. Context should be represented in a form the system can inspect, restrictions should be applied before relevance scoring, and evaluation should measure several competing outcomes. Holistic Mind uses short structured questions to make the first principle concrete, shared comfort mapping for the second, and separate quality, fill, and violation measures for the third. Together, these choices make failures easier to locate than a single opaque model output would permit.

There is also a clear gap between recommendation research and this application's eventual use. General-purpose semantic models are trained for broad language similarity rather than this exercise catalogue. Authored relevance judgments may reward descriptions that match familiar language without establishing personal usefulness. The project therefore needs domain-specific review and user-centred evidence before moving from technical feasibility to stronger effectiveness claims. This gap motivates the critical evaluation rather than undermining the value of building and testing a bounded prototype.

\subsection{Trauma-informed principles as design requirements}

SAMHSA describes trauma-informed approaches through principles including safety, transparency, peer support, collaboration, empowerment and choice, and attention to cultural and historical context \citep{samhsa}. These principles originate in wider organisational practice. Applying them to a mobile interface is an engineering interpretation that requires evaluation with its intended users. Holistic Mind chiefly operationalises choice, transparency, and predictable interaction; it does not implement every part of a comprehensive trauma-informed service.

Choice appears in optional practice selection and explicit comfort preferences. Predictability appears in short onboarding, visible navigation, and recognisable controls. Transparency appears in recommendation reasons and the explanation that journal recovery requires a separate key. The design avoids requiring a trauma narrative as the price of receiving a suggestion. Personalisation uses structured self-reports rather than an attempt to classify a person's history. The relevant question is whether these choices feel useful and respectful in practice, which cannot be answered by a software test alone.

Peer support is a clear example of a principle outside the present application. There is no evidenced moderated peer community or therapist relationship inside Holistic Mind. Cultural responsiveness is also incomplete: accepting Nepali text in an encrypted journal demonstrates text handling, not a fully localised or culturally validated intervention. Treating partial implementation honestly makes the term trauma-informed more precise. It identifies a design direction and review framework rather than certifying the application's overall suitability.

\subsection{Digital mental health evidence and its limits}

\citet{linardon2019} reviewed 66 randomised trials of app-supported smartphone interventions. Their findings support the potential of some interventions for several mental health outcomes, while noting variable trial quality and limited evidence for some outcomes. The review does not establish that Holistic Mind's specific content, recommendation rules, or intended audience will benefit. It also does not substantiate the earlier report's claim that fewer than 3\% of apps use trauma frameworks. That percentage is consequently excluded from the final argument.

The distinction matters because a functioning application and an effective intervention are different achievements. Authentication, persistent journaling, and a reliable media player can be evaluated through technical and usability methods. Claims about symptom improvement require an appropriate study design and outcome measures. Holistic Mind's evaluation is presently concentrated on software functionality and relevance agreement with authored labels. The literature motivates further research, but its findings are not transferred as outcomes of this artefact.

\subsection{Related recommender platforms and corrected sources}

The earlier interim report identified work on recommending mental health applications. Its linked 2024 paper is authored by Tapuria and colleagues, rather than the author attribution used in that draft. \citet{tapuria2024} describe requirements assessment, prototype development, and bench testing for a platform that recommends trusted mental health apps. Their emphasis on safety, confidentiality, and reliability is relevant to Holistic Mind, although recommending complete apps differs from ranking practices within one catalogue.

This comparison reinforces the importance of content quality and user requirements before algorithm complexity. Holistic Mind uses administered exercise metadata rather than choosing from independently assessed third-party applications. The administrator's publication decision therefore becomes part of the recommendation environment. Similarity scores cannot compensate for missing restrictions or unsupported descriptions. The project needs a content-review procedure alongside ranking evaluation, especially when the catalogue expands or descriptions change.

The NarraGive evaluation offers another relevant comparison. It examines content-based and collaborative recommendation of mental health recovery narratives, using measures including accuracy, precision, diversity, coverage, and unfairness \citep{slade2024}. It also respects user-blocked narratives. Holistic Mind uses different content and a much smaller authored evaluation, but the shared lesson is to assess more than relevance alone. The paper's participant outcomes and dataset are not evidence about Holistic Mind.

\citet{matthews2025} provide a scoping review specifically addressing personalisation and recommendation for mental health apps. This is a stronger source description than the abbreviated title and author details in the earlier interim bibliography. It supports positioning the project within an established research topic rather than presenting personalisation as an entirely new idea. The final report uses verified source identities and avoids inheriting unverified claims of market uniqueness.
\#\# 2.9 Literature synthesis and project positioning

Four bodies of work shape the final artefact: trauma-informed principles, evidence on digital interventions, hybrid recommendation, and privacy research. Their roles differ. Trauma-informed principles guide interaction decisions. Digital intervention evidence motivates research without validating this application. Hybrid recommendation provides a practical way to combine heterogeneous evidence. Privacy research challenges the assumption that semantic representations adequately conceal personal text. Keeping these roles distinct prevents a technology choice from becoming an unsupported clinical argument.

The project is positioned as an applied software contribution. Its originality lies in assembling an explicit daily workflow, comfort-aware eligibility, inspectable ranking, managed learning content, and device-side reflection within one system. It does not claim a first-of-its-kind app, a novel neural architecture, or an exhaustive market gap. The original proposal's focus remains meaningful without asserting that no competing application addresses similar needs. The stronger academic argument is to show the implemented mechanisms, evaluate them fairly, and identify the limits of the resulting evidence.
